\chapter{Database CRUD with Records}
\label{ch:db-crud}

The \texttt{Record}\index{Record} class represents a single database row.  This chapter covers creating, reading, updating, and deleting records using the Kiss pattern, which replaces raw SQL for most single-table operations.

\section{The Record Class}
\label{sec:crud-record}

A \texttt{Record} holds column name/value pairs for one row.  Records come from two sources:

\begin{enumerate}
  \item \textbf{Query results} --- \texttt{fetchOne()}, \texttt{fetchAll()}, and cursor operations return Records populated with data from the database.  These records know their table name and primary key, so they can be updated or deleted.
  \item \textbf{New records} --- \texttt{db.newRecord("table")} creates an empty Record for a specified table.  You populate it with \texttt{set()} calls and then insert it with \texttt{addRecord()}.
\end{enumerate}

\section{Creating and Inserting Records}
\label{sec:crud-insert}

The Kiss pattern for inserting a row:

\begin{lstlisting}[style=groovy]
Record rec = db.newRecord("phone")
rec.set("first_name", "Alice")
rec.set("last_name", "Smith")
rec.set("phone_number", "555-1234")
rec.addRecord()
\end{lstlisting}

\noindent
Only columns that are explicitly set are included in the INSERT statement.  Columns with database defaults (auto-increment, \texttt{CURRENT\_TIMESTAMP}, etc.) are left out and receive their default values.

\begin{notebox}
Always use the \texttt{newRecord}/\texttt{set}/\texttt{addRecord} pattern rather than raw SQL \texttt{INSERT} statements.  It is cleaner, handles column types automatically, and only includes non-null columns in the INSERT.
\end{notebox}

\subsection{Auto-Increment Keys}

If the table has an auto-increment primary key and you need the generated value:

\begin{lstlisting}[style=groovy]
Record rec = db.newRecord("employee")
rec.set("first_name", "Bob")
rec.set("last_name", "Jones")
Object generatedId = rec.addRecordAutoInc()
int employeeId = (int) generatedId
\end{lstlisting}

\noindent
\texttt{addRecordAutoInc()} inserts the record and returns the generated primary key value.  Cast it to the appropriate type (\texttt{int}, \texttt{long}, or \texttt{short}).

\subsection{Method Chaining}

The \texttt{set()} methods return the Record, allowing method chaining:

\begin{lstlisting}[style=groovy]
db.newRecord("phone")
    .set("first_name", "Alice")
    .set("last_name", "Smith")
    .set("phone_number", "555-1234")
    .addRecord()
\end{lstlisting}

\section{Fetching Records}
\label{sec:crud-fetch}

\begin{lstlisting}[style=groovy]
// Single record by primary key
Record emp = db.fetchOne(
    "SELECT * FROM employee WHERE employee_id = ?", empId)

// Multiple records
List<Record> active = db.fetchAll(
    "SELECT * FROM employee WHERE active = 'Y' " +
    "ORDER BY last_name")
\end{lstlisting}

\noindent
\texttt{fetchOne()} returns \texttt{null} if no rows match.  \texttt{fetchAll()} returns an empty list if no rows match.

\begin{warnbox}
For a Record to be updatable or deletable, the SELECT must:
\begin{itemize}
  \item Query a single table (no joins).
  \item Include the primary key column(s) in the result.
\end{itemize}
\end{warnbox}

\section{Reading Field Values}
\label{sec:crud-read}

Records provide typed getter methods for each data type:

\begin{lstlisting}[style=groovy]
String name       = rec.getString("first_name")
int id            = rec.getInt("employee_id")
double salary     = rec.getDouble("salary")
long bigNum       = rec.getLong("big_number")
short small       = rec.getShort("small_value")
float pct         = rec.getFloat("percentage")
BigDecimal amount = rec.getBigDecimal("invoice_total")
boolean flag      = rec.getBoolean("active_flag")
char code         = rec.getChar("status_code")
byte[] data       = rec.getByteArray("binary_data")
\end{lstlisting}

\noindent
All getters return \texttt{null} if the column value is NULL.  Column names are case-insensitive.

\noindent
\texttt{getBigDecimal()} is the right choice for \texttt{NUMERIC} and \texttt{DECIMAL} columns where exact precision matters (money, tax rates, quantities with fixed scale).  Many JDBC drivers return these columns as \texttt{BigDecimal} natively, so \texttt{getInt()}, \texttt{getLong()}, \texttt{getFloat()}, and \texttt{getDouble()} also accept a \texttt{BigDecimal}-backed column and convert it --- exactly for the integer variants (throwing if the value does not fit), and with possible loss of precision for \texttt{float}/\texttt{double}.

\subsection{Date and Time Fields}

Kiss provides several methods for reading date and time values:

\begin{lstlisting}[style=groovy]
// Date with time information
Date hired = rec.getDateTime("hire_date")

// Date portion only (time zeroed out)
Date dateOnly = rec.getDateOnly("birth_date")

// Date as YYYYMMDD integer
int dateInt = rec.getDateAsInt("birth_date")  // e.g. 19950315

// Date/time as milliseconds since epoch
long ms = rec.getDateTimeMS("created_at")

// Time portion only (milliseconds)
long time = rec.getTime("start_time")

// ZonedDateTime (preserves timezone)
ZonedDateTime zdt = rec.getZonedDateTime("event_time")
\end{lstlisting}

\begin{warnbox}
\texttt{getDateTime()} returns \texttt{java.util.Date}, \emph{not} \texttt{java.sql.Timestamp}.  Never cast the result to \texttt{Timestamp}.  This is a common mistake that causes \texttt{ClassCastException} at runtime.
\end{warnbox}

\subsection{Setting Date and Time Fields}

\begin{lstlisting}[style=groovy]
// Date and time
rec.setDateTime("created_at", new Date())

// Date portion only (no time)
rec.setDateOnly("birth_date", new Date())

// Date from YYYYMMDD integer
rec.setDateOnly("birth_date", 19950315)

// Date/time from milliseconds since epoch
rec.setDateTime("event_time", System.currentTimeMillis())

// Time from milliseconds
rec.setTime("start_time", System.currentTimeMillis())
\end{lstlisting}

\section{Updating Records}
\label{sec:crud-update}

Fetch the record, change the fields, and call \texttt{update()}:

\begin{lstlisting}[style=groovy]
Record emp = db.fetchOne(
    "SELECT * FROM employee WHERE employee_id = ?", id)
emp.set("salary", newSalary)
emp.set("active", "Y")
emp.update()
\end{lstlisting}

\noindent
\texttt{update()} is smart: it compares the current column values against the original values fetched from the database and only includes changed columns in the UPDATE statement.  If nothing has changed, no SQL is executed.

\begin{tipbox}
Because \texttt{update()} only sends changed fields, it is safe to fetch a record, modify one column, and call \texttt{update()} without worrying about overwriting concurrent changes to other columns.
\end{tipbox}

\section{Deleting Records}
\label{sec:crud-delete}

\begin{lstlisting}[style=groovy]
Record emp = db.fetchOne(
    "SELECT * FROM employee WHERE employee_id = ?", id)
emp.delete()
\end{lstlisting}

\noindent
Alternatively, use \texttt{Connection.execute()} for a direct delete without fetching:

\begin{lstlisting}[style=groovy]
db.execute("DELETE FROM employee WHERE employee_id = ?", id)
\end{lstlisting}

\subsection{Delete Callbacks}

You can register a connection-wide callback that fires whenever a record is deleted:

\begin{lstlisting}[style=groovy]
db.setDeleteCallback({ String table, Object pkValue ->
    logger.info("Deleted from ${table}: ${pkValue}")
})
\end{lstlisting}

\noindent
\texttt{ExternalFile} uses this exact mechanism to remove a deleted row's associated file(s) automatically, opted in one table at a time --- see Section~\ref{sec:file-external-cascade}.

\section{Checking Column Existence}
\label{sec:crud-colexist}

\begin{lstlisting}[style=groovy]
if (rec.columnExists("middle_name")) {
    String middle = rec.getString("middle_name")
}
\end{lstlisting}

\section{Copying Records}
\label{sec:crud-copy}

\begin{lstlisting}[style=groovy]
// Copy all columns from one record to another
newRec.copy(existingRec)

// Copy only columns that exist in both records (by name
// and type)
newRec.copyCorresponding(existingRec)
\end{lstlisting}

\noindent
\texttt{copyCorresponding()} is useful when copying between records from different tables that share some column names.

\section{Converting to JSON}
\label{sec:crud-json}

Records can be converted to JSON for direct frontend communication:

\begin{lstlisting}[style=groovy]
// Single record to JSON object
JSONObject json = emp.toJSON()
outjson.put("employee", json)

// Add record columns to an existing JSON object
emp.addToJSON(outjson)

// Convert a list of records to a JSON array
List<Record> recs = db.fetchAll("SELECT * FROM employee")
JSONArray arr = Record.toJSONArray(recs)
outjson.put("employees", arr)
\end{lstlisting}

\noindent
For the common case, use \texttt{fetchAllJSON()} or \texttt{fetchOneJSON()} directly on the Connection --- these skip the intermediate Record objects:

\begin{lstlisting}[style=groovy]
// Query directly to JSON array
outjson.put("records", db.fetchAllJSON(
    "SELECT employee_id, first_name, last_name FROM employee"))

// Query directly to JSON object
outjson.put("employee", db.fetchOneJSON(
    "SELECT * FROM employee WHERE employee_id = ?", id))
\end{lstlisting}

\section{A Complete CRUD Service}
\label{sec:crud-example}

The following example shows a complete service with list, add, edit, and delete methods, along with the corresponding frontend screen.

\subsection{Backend Service}

\begin{lstlisting}[style=groovy]
package services

class EmployeeService {

    void list(JSONObject injson, JSONObject outjson,
              Connection db, ProcessServlet servlet) {
        outjson.put("employees", db.fetchAllJSON(
            "SELECT employee_id, first_name, last_name, " +
            "active FROM employee " +
            "ORDER BY last_name, first_name"))
    }

    void get(JSONObject injson, JSONObject outjson,
             Connection db, ProcessServlet servlet) {
        Record emp = db.fetchOne(
            "SELECT * FROM employee WHERE employee_id = ?",
            injson.getInt("employeeId"))
        if (emp == null)
            throw new UserException("Employee not found")
        emp.addToJSON(outjson)
    }

    void create(JSONObject injson, JSONObject outjson,
                Connection db, ProcessServlet servlet) {
        Record rec = db.newRecord("employee")
        rec.set("first_name", injson.getString("firstName"))
        rec.set("last_name", injson.getString("lastName"))
        rec.set("active", "Y")
        rec.addRecord()
    }

    void save(JSONObject injson, JSONObject outjson,
              Connection db, ProcessServlet servlet) {
        Record rec = db.fetchOne(
            "SELECT * FROM employee WHERE employee_id = ?",
            injson.getInt("employeeId"))
        rec.set("first_name", injson.getString("firstName"))
        rec.set("last_name", injson.getString("lastName"))
        rec.set("active", injson.getString("active"))
        rec.update()
    }

    void delete(JSONObject injson, JSONObject outjson,
                Connection db, ProcessServlet servlet) {
        db.execute(
            "DELETE FROM employee WHERE employee_id = ?",
            injson.getInt("employeeId"))
    }
}
\end{lstlisting}

\subsection{Frontend Screen (HTML)}

\begin{lstlisting}[style=html]
<div>
  <h3>Employees</h3>
  <div id="emp-grid" style="width: 100%; height: 300px;"></div>
  <div style="margin-top: 10px;">
    <push-button id="emp-new">New</push-button>
    <push-button id="emp-edit" disabled>Edit</push-button>
    <push-button id="emp-delete" disabled>Delete</push-button>
  </div>

  <popup id="emp-popup" height="160px" width="400px">
    <popup-title>Employee</popup-title>
    <popup-body>
      <div style="display: grid;
                  grid-template-columns: auto 1fr;
                  gap: 8px; padding: 10px;">
        <span>First Name:</span>
        <text-input id="ep-first" required
                    maxlength="50"></text-input>
        <span>Last Name:</span>
        <text-input id="ep-last" required
                    maxlength="50"></text-input>
      </div>
      <div style="text-align: right; margin-top: 10px;">
        <push-button id="ep-ok">OK</push-button>
        <push-button id="ep-cancel">Cancel</push-button>
      </div>
    </popup-body>
  </popup>
</div>
\end{lstlisting}

\subsection{Frontend Screen (JavaScript)}

\begin{lstlisting}[style=javascript]
'use strict';

(async function () {
    const WS = 'services/EmployeeService';

    const grid = new AGGrid('emp-grid', [
        {headerName: 'Last Name', field: 'last_name',
            width: 40},
        {headerName: 'First Name', field: 'first_name',
            width: 40},
        {headerName: 'Active', field: 'active', width: 20}
    ], 'employee_id');
    grid.show();

    async function loadGrid() {
        grid.clear();
        $$('emp-edit').disable();
        $$('emp-delete').disable();
        const res = await Server.call(WS, 'list');
        if (res._Success)
            grid.addRecords(res.employees);
    }

    loadGrid();

    grid.setOnSelectionChanged((rows) => {
        $$('emp-edit').enable(rows);
        $$('emp-delete').enable(rows);
    });

    $$('emp-new').onclick(() => {
        $$('ep-first').clear();
        $$('ep-last').clear();
        Utils.popup_open('emp-popup', 'ep-first');

        $$('ep-ok').onclick(async () => {
            if ($$('ep-first').isError('First Name')) return;
            if ($$('ep-last').isError('Last Name')) return;
            const res = await Server.call(WS, 'create', {
                firstName: $$('ep-first').getValue(),
                lastName: $$('ep-last').getValue()
            });
            if (res._Success) {
                Utils.popup_close();
                loadGrid();
            }
        });

        $$('ep-cancel').onclick(() => Utils.popup_close());
    });

    $$('emp-delete').onclick(() => {
        Utils.yesNo('Confirm', 'Delete this employee?',
                    async () => {
            const row = grid.getSelectedRow();
            const res = await Server.call(WS, 'delete', {
                employeeId: row.employee_id
            });
            if (res._Success) loadGrid();
        });
    });
})();
\end{lstlisting}

\noindent
This example demonstrates the standard Kiss full-stack pattern: a backend service with CRUD methods, a frontend screen with a grid and popup editor, and \texttt{Server.call()} connecting the two.
